\section{Inhomogeneous Ising chains}

Throughout, $S = \mathbb N$, $E = \set{\pm 1}$ with the uniform a priori measure, $J_n > 0$ the
couplings, $\Lambda_N = \set{0, \dots, N-1}$, and $\gamma$ the Gibbs specification of the
nearest-neighbour potential $\Phi_{\set{n, n+1}}(\sigma) = -J_n \sigma_n \sigma_{n+1}$ (6.2).

\begin{lemma}[Georgii (6.5)--(6.6)]
    \label{lem:iic-6-5}
    \uses{def:pot-gibbs-spec, def:sb-pair-potential}
    \lean{MeasureTheory.GibbsMeasure.isingChainSpecification, MeasureTheory.GibbsMeasure.chainZ_eq_prod, MeasureTheory.GibbsMeasure.chainMag_eq, MeasureTheory.GibbsMeasure.isingChainSpecification_range_apply_setOf_eq, MeasureTheory.GibbsMeasure.isingChainSpecification_range_apply_eq_sum}
    \leanok

    $Z_{\Lambda_N} = 2^N \prod_{i < N} \cosh J_i$, and for $n < N$
    \[
        \gamma_{\Lambda_N}(\sigma_n = x \mid \omega)
        = \tfrac12\Bigl(1 + x\,\omega_N \prod_{i=n}^{N-1} \tanh J_i\Bigr) ; \tag{6.5}
    \]
    for $A \in \mathcal F_{\Lambda_n}$ and $n \leq N$,
    $\gamma_{\Lambda_N}(A \mid \omega) = \sum_x \gamma_{\Lambda_n}(A \mid x)\,\gamma_{\Lambda_N}(\sigma_n = x \mid \omega)$ (6.6).
\end{lemma}
\begin{proof}
    \leanok
    Transfer-matrix recursion in the boundary spin; (6.6) is the Markov property of $\gamma$
    composed with the consistency $\gamma_{\Lambda_N}\gamma_{\Lambda_n} = \gamma_{\Lambda_N}$.
\end{proof}

\begin{theorem}[Georgii (6.4)]
    \label{thm:iic-6-4}
    \uses{lem:iic-6-5, def:gibbs-meas, thm:exg-4-17}
    \lean{MeasureTheory.GibbsMeasure.exists_mem_GP_tendsto, MeasureTheory.GibbsMeasure.eq_smul_add_smul_of_isGibbsMeasure, MeasureTheory.GibbsMeasure.exists_G_isingChainSpecification_eq, MeasureTheory.GibbsMeasure.exists_mem_G_integral_spin_pos}
    \leanok

    Suppose $\sum_n e^{-2J_n} < \infty$ (6.1). Then there are $\mu_+, \mu_- \in \mathcal G(\gamma)$
    with $\mu_- = \tau\mu_+$ for the spin flip $\tau$, $\mu_+(\sigma_n) > 0 > \mu_-(\sigma_n)$
    for all $n$, and $\mathcal G(\gamma) = \set{s\mu_+ + (1-s)\mu_- : 0 \leq s \leq 1}$; in
    particular $|\mathcal G(\gamma)| > 1$.
\end{theorem}
\begin{proof}
    \leanok
    (6.1) makes $\prod_{i \geq n}\tanh J_i$ converge to a positive limit, so by (6.6)
    $\gamma_{\Lambda_N}(\cdot \mid \omega)$ converges on local events for every $\omega$ eventually
    $+1$; a cluster point (Theorem (4.17)) is the limit $\mu_+$, and $\mu_- = \tau\mu_+$ is the
    limit for $\omega$ eventually $-1$. For $\mu \in \mathcal G(\gamma)$, $\mu(\sigma_n \neq \sigma_{n+1}) = \frac12(1 - \tanh J_n) \leq e^{-2J_n}$,
    so by Borel--Cantelli $\mu(A_+ \cup A_-) = 1$ with $A_\pm$ the eventually-constant
    configurations, and the DLR equations give $\mu = \mu(A_+)\mu_+ + \mu(A_-)\mu_-$.
\end{proof}

\begin{remark}[Georgii (6.7)(2)]
    \label{rmk:iic-6-7}
    \uses{thm:iic-6-4}
    \lean{MeasureTheory.GibbsMeasure.exists_G_eq_singleton_of_not_summable}
    \leanok

    If $\sum_n e^{-2J_n} = \infty$ then $\prod_{i \geq n} \tanh J_i = 0$, the limit in (6.6) is
    independent of $\omega$, and $|\mathcal G(\gamma)| = 1$.
\end{remark}

\begin{remark}[Georgii (6.7)(1), the low-temperature limit]
    \label{rmk:iic-6-7-1}
    \uses{thm:iic-6-4, def:lc-topology}
    \lean{MeasureTheory.GibbsMeasure.plusPhase, MeasureTheory.GibbsMeasure.minusPhase, MeasureTheory.GibbsMeasure.exists_G_gibbsSpecification_smul_eq, MeasureTheory.GibbsMeasure.integral_spin_plusPhase, MeasureTheory.GibbsMeasure.tendsto_integral_spin_plusPhase_smul_atTop, MeasureTheory.GibbsMeasure.tendsto_localDist_plusPhase_smul_atTop, MeasureTheory.GibbsMeasure.tendsto_localDistSet_GP_gibbsSpecification_dirac, MeasureTheory.GibbsMeasure.exists_eq_smul_dirac_add_smul_dirac_of_tendsto}
    \leanok

    Let $J_n > 0$ with $\sum_n e^{-2J_n} < \infty$. The constant configurations $\omega^\pm$ are
    ground states, $\mathcal G(\beta\Phi) = [\mu_-^\beta, \mu_+^\beta]$ for $\beta \geq 1$,
    $\mu_+^\beta(\sigma_n) = \prod_{i \geq n}\tanh(\beta J_i) \to 1$ as $\beta \to \infty$, and
    $\mu_\pm^\beta \to \delta_{\omega^\pm}$ in the topology of local convergence; every local
    limit of $\mu^\beta \in \mathcal G(\beta\Phi)$ is of the form
    $s\delta_{\omega^+} + (1-s)\delta_{\omega^-}$.
\end{remark}
\begin{proof}
    \leanok
    $\mu_\pm(\sigma_a \neq \pm 1) \leq \sum_i e^{-2J_i}$ uniformly in $a$, and
    $\sum_i e^{-2\beta J_i} \to 0$ by Tannery's theorem.
\end{proof}

\begin{remark}[Georgii (6.7)(3), percolation]
    \label{rmk:iic-6-7-3}
    \uses{thm:iic-6-4, thm:ext-77-concrete}
    \lean{MeasureTheory.GibbsMeasure.plusPhase_apply_eventuallyConst_true, MeasureTheory.GibbsMeasure.plusPhase_mem_extremePoints_G, MeasureTheory.GibbsMeasure.minusPhase_mem_extremePoints_G, MeasureTheory.GibbsMeasure.chainGraph, MeasureTheory.GibbsMeasure.chainGraph_adj_iff_forall_le, MeasureTheory.GibbsMeasure.mem_eventuallyConst_union_iff_existsUnique_infinite_supp, MeasureTheory.GibbsMeasure.apply_setOf_exists_infinite_supp_eq_one, SimpleGraph.exists_infinite_supp_iff}
    \leanok

    Under (6.1), $\mu_+(A_+) = 1 = \mu_-(A_-)$ for the disjoint tail events
    $A_\pm = \{\omega \text{ eventually } \pm 1\}$, and $\mu_\pm$ are the extreme points of
    $\mathcal G(\gamma)$. For $\omega$ let $G(\omega)$ be the graph on $\mathbb N$ with a bond
    $\{i, i+1\}$ iff $\omega_i = \omega_{i+1}$, i.e.\ iff $\Phi_{\{i,i+1\}}(\omega)$ is minimal.
    Then $\omega \in A_+ \cup A_-$ iff $G(\omega)$ has a (unique) infinite connected component,
    which has $\mu$-probability one for every $\mu \in \mathcal G(\gamma)$.
\end{remark}
\begin{proof}
    \leanok
    A subgraph of the half-line has an infinite component iff all but finitely many bonds are
    present; $\mu_\pm(A_\pm) = 1$ from the site probabilities of (6.4) and Borel--Cantelli.
\end{proof}
